% Needs \usepackage{natbib} (for \citet) and this entry in the .bib file:
% @mastersthesis{baggenstos2026,
%   author = {Baggenstos, Matteo},
%   title  = {Sequential {MATSim}--{SUMO} Hybrid Simulation: A Case Study in {Zurich} ({Switzerland})},
%   school = {IVT, ETH Z{\"u}rich},
%   year   = {2026},
%   month  = jun,
%   type   = {Master's thesis}
% }
% and for the signal timing and the count check:
% @techreport{webster1958,
%   author = {Webster, F. V.}, title = {Traffic Signal Settings},
%   institution = {Road Research Laboratory, HMSO}, address = {London},
%   type = {Road Research Technical Paper}, number = {39}, year = {1958}
% }
% @techreport{koonce2008,
%   author = {Koonce, Peter and Rodegerdts, Lee and others}, title = {Traffic Signal Timing Manual},
%   institution = {Federal Highway Administration}, number = {FHWA-HOP-08-024}, year = {2008}
% }
% @manual{dmrb1996,
%   title = {Design Manual for Roads and Bridges, Volume 12, Section 2, Part 1: Traffic Appraisal in Urban Areas},
%   organization = {Highways Agency}, address = {London}, year = {1996}
% }
% @misc{zurichcounts,
%   author = {{Stadt Z{\"u}rich, Dienstabteilung Verkehr}},
%   title = {Daten der Verkehrsz{\"a}hlung zum motorisierten Individualverkehr (Stundenwerte)},
%   howpublished = {Open Data Z{\"u}rich, dataset sid\_dav\_verkehrszaehlung\_miv\_od2031, CC0}, year = {2024}
% }

% ===== (B) replaces \subsubsection{Evolutionary Repair of Remaining Gridlocks} =====
\subsubsection{Evolutionary Repair of Remaining Gridlocks}
\label{subsubsec:evolutionary-turns}

The filtering rules repair what the network shows on its own. Some
errors only appear when cars use the network: a lane that leads to a
junction its cars cannot cross, a missing turn, or a lane-change ban
that traps a car in the wrong lane. This stage finds them by simulating
the National demand on the corrected City of Zürich network and
repairs them with an evolutionary search. It reads the corrected file
and writes each new network beside it. Figure~\ref{fig:repair-flow}
shows the procedure.

The demand is the eqasim National car and truck demand for 07:00 to
08:00, cut to the Zürich network as described in
Section~\ref{subsec:demand-transfer}. Trips are routed with
\texttt{duarouter}, and the hour is simulated in SUMO. A vehicle that
has not moved for \(120\,\mathrm{s}\) is removed and counted. SUMO
records why it waited: in the wrong lane, at a lane with no outgoing
connection, behind a queue, while yielding, or after a collision. A
trip with no path through the network is counted as unroutable.

The repair runs in rounds, and each round has three steps.

\paragraph{OpenStreetMap step.} The network is first compared with
OpenStreetMap. A link that ends at a junction where the OpenStreetMap
way continues, but has no SUMO connection onward, is connected to that
continuation. Links that join otherwise separate parts of the drivable
network are connected in the same way. The step is repeated up to four
times, because a new connection can expose the next missing one.

\paragraph{Simulation and candidate edits.} The National demand is
simulated on the checked network. The edits that the search may make
come from the places where vehicles stopped. To bound the number of
Google requests, only the 40 approaches with the most stuck vehicles
are used. Three kinds of edit are formed.

\begin{enumerate}[label=\textbf{G\arabic*:}]
    \item \emph{Turns.} At each approach, the usual left, straight, and
    right movements are submitted to the Google Routes API as driving
    routes with traffic ignored. The origin lies \(40\,\mathrm{m}\)
    upstream of the junction and the destination \(40\,\mathrm{m}\)
    downstream. A route no longer than \(\max(2.5\,d,\ d+250\,\mathrm{m})\),
    where \(d\) is the direct distance through the junction, is read as
    a movement that exists. A longer route is read as a movement Google
    does not drive. An edit is formed only where Google disagrees with
    the network: a missing movement that Google drives may be added,
    and a present movement that Google does not drive may be deleted.

    \item \emph{Lane movements.} For a vehicle that waited in the wrong
    lane, its route is followed from that lane through the lane
    connections, for up to \(1{,}000\,\mathrm{m}\). Where a lane-change
    ban keeps the vehicle in its lane, it stays there. The first
    junction at which the lane cannot reach the next edge of the route
    gives the blocked movement. When Google drives that movement, a
    connection from that lane is a candidate edit. At most twelve such
    movements are used, starting with those that stopped the most
    vehicles.

    \item \emph{Lane-change bans.} SUMO imports the OpenStreetMap
    \texttt{change:lanes} tags as bans on changing lanes. A ban on the
    approach to a junction can keep a car in a lane that does not lead
    where it is going. For each edge with wrong-lane stops, the edges up
    to \(500\,\mathrm{m}\) downstream that carry a ban are collected,
    and each connected stretch of them is one candidate: lifting the
    ban on that stretch.
\end{enumerate}

Each candidate is one gene, and a chromosome is a binary vector that
switches the candidates on or off. A chromosome is turned into a
network with \texttt{netconvert}, which adds or deletes the chosen
connections, and the chosen bans are then lifted. Routes are repaired
for the new network before it is simulated. The score of a network is
\begin{equation}
    F = n_{\mathrm{stuck}} + n_{\mathrm{unroutable}} + 0.001\,n_{\mathrm{edits}},
    \label{eq:repair-fitness}
\end{equation}
where \(n_{\mathrm{stuck}}\) is the number of stuck vehicles,
\(n_{\mathrm{unroutable}}\) the number of trips without a path, and
\(n_{\mathrm{edits}}\) the number of genes switched on. Counting
unroutable trips means that deleting a movement cannot improve the
score by removing the cars that used it. The small edit term prefers
the smaller of two networks with the same count.

\paragraph{Evolutionary search.} The population holds eight
chromosomes, one of them the network without further edits. The two
best individuals are kept in each generation. The others are produced
by uniform crossover and by flipping one gene. The search stops when
the best score reaches zero, when two successive generations do not
improve it, or after six generations.

Every individual is compared with OpenStreetMap after it is simulated.
It is rejected when it has a link to nowhere that the network at the
start of the round did not have, or when it splits the drivable network
into more pieces. A rejected individual cannot be selected, whatever
its score.

The best accepted network starts the next round. The rounds stop when
no vehicle is stuck, when the search keeps no edit or does not improve
on the start of the round, or after the maximum number of rounds. The
last network then receives a final OpenStreetMap check and a final
simulation, and the places where vehicles still stop are listed.

\paragraph{Sampled and full demand.} The eqasim scenario simulates a
sample of the population, so the repair is run in stages. Stage~1 uses
the sampled demand for up to ten rounds. Errors in the layout of the
network stop cars whatever the volume, so this stage finds most of
them, and each simulation takes about a minute. Stage~2 scales the
demand to the full population as described in
Section~\ref{subsec:demand-scaling} and runs up to three further rounds
from the network that stage~1 produced. The unrepaired network is
simulated with the same demand for comparison. From stage~2 on, a
vehicle counts as stuck only after \(300\,\mathrm{s}\) without moving,
the SUMO default, because waiting two minutes at a busy junction in the
morning peak is normal. Stuck vehicles are counted from SUMO's
statistic output (\texttt{--statistic-output}): the teleports per
inserted vehicle, split into queue, yield, wrong lane, and collision,
and the vehicles that could not enter the network before the hour
ended. A collision is resolved by one teleport, the SUMO default; with
collisions only reported, SUMO repeats the warning every step, so
counting warnings would count one collision many times. At full volume
some vehicles stop because a junction or a road has too little
capacity. Turn and lane edits cannot add capacity, so the remaining
stops are reported in two groups: network errors (wrong lane, no
connection, and collision) and capacity (queue and yield).

\paragraph{Routes and lane-change bans.} From stage~2 on, the trips are
routed once with \texttt{duarouter}, and every later simulation uses
these routes and the same SUMO options, so a result can be repeated in
\texttt{sumo-gui}. Each vehicle also chooses its route again at
departure and every five minutes from the current travel times (SUMO
rerouting device), instead of all clones of a trip driving the
shortest path of the empty network. Both routers choose a sequence of
edges and assume that a vehicle can change to any lane on the way.
Where OpenStreetMap marks a lane change as forbidden
(\texttt{change:lanes}, solid lines), some of these routes cannot be
driven, and the vehicle waits at the end of its lane until it is
teleported. At the A1/A4 merge at Opfikon, for example, the lanes to
Zürich-City are separated from the lanes to St.~Gallen by a solid line,
so a car from the Zürich-City slip road cannot continue towards
St.~Gallen. Every route is therefore checked lane by lane, using only
the network's connections and the lane changes allowed for cars. A
route that cannot be driven is replaced by the fastest route that can
(free-flow times, Dijkstra on the lane graph). A vehicle whose route
passes an edge from which an undrivable path starts within
\(1\,\mathrm{km}\) keeps its route and is not rerouted.

\paragraph{Signal timing.} The traffic lights of the network are not
the programs of the city: netconvert guessed which junctions have
signals from OpenStreetMap and gave each an actuated program with
standard phase lengths. Before the behaviour calibration, the cycle
length and the green splits of every signal are set from the routed
full demand with Webster's method \citep{webster1958}, using SUMO's
\texttt{tlsCycleAdaptation.py} (saturation headway \(2\,\mathrm{s}\),
lost time \(4\,\mathrm{s}\) per phase, cycle
\(60\)--\(120\,\mathrm{s}\), a truck counting as 3.5 cars). Yellow
and all-red phases keep their lengths. The signals stay actuated, as
most signals in Zürich are traffic-actuated: a green phase starts with
Webster's green, keeps its minimum green, and may be extended by the
detectors up to 1.25 times Webster's green, the lower end of the usual
maximum green of an actuated controller \citep{koonce2008}. A detector
that is occupied for \(30\,\mathrm{s}\) counts as jammed, so that
SUMO also uses detectors on lanes shared by movements with different
phases. For comparison, the guessed programs and a fixed-time Webster
plan with offsets for green waves (SUMO's \texttt{tlsCoordinator.py})
are simulated as well; the actuated plan is used in any case, so the
choice does not depend on the number of stuck vehicles. Where an edge
shorter than one car (\(7.5\,\mathrm{m}\)) separates two junctions and
at least two pairs of vehicles collided on it, the two junctions are
joined into one, since a car stopped on such an edge stands across the
junction behind it.

\paragraph{Driver behaviour.} Stage~3 calibrates how drivers follow,
yield, and change lanes, on the network that stage~2 produced and with
the Webster signal programs. It runs after the network repair so that
a behaviour setting cannot hide a lane or connection error. The
network, its lanes, and the signal programs are not changed, and no
parameter lets a driver ignore a red light.
Table~\ref{tab:behaviour-genes} lists the eleven genes and their
ranges; each range contains the SUMO default. The genes are real
numbers. The search uses the same scheme as the network repair: six
individuals, two elites, tournament selection, uniform crossover, and
one gene redrawn per child, for three generations. The first population
contains the SUMO defaults and the microscopic settings of
\citet{baggenstos2026} (\texttt{ignore-junction-blocker}
\(1\,\mathrm{s}\), \texttt{time-to-impatience} \(30\,\mathrm{s}\)).

Removing every queue would make the model faster than the real city.
The score therefore combines the stuck vehicles with the travel time.
Sixty trips of the scaled demand that depart in the first half hour are
drawn once. Their travel times are requested from the Google Routes API
with traffic, for the same departure time on a weekday, as
\citet{baggenstos2026} did with TomTom. With \(p\) the teleported share
in percent, \(w\) the share of the loaded vehicles that could not enter
the network before the hour ended in percent, and \(b\) the bias of the
mean SUMO travel time against Google in percent, the score is
\begin{equation}
    F_{\mathrm{b}} = p + w + \max\bigl(0,\ |b| - 10\bigr).
    \label{eq:behaviour-fitness}
\end{equation}
Travel times within \(\pm 10\,\%\) of Google are not penalised; beyond
that, each percentage point counts as much as one percent of teleported
vehicles. Without \(w\), a setting that blocks the entries of the
network would score well, because a vehicle that never enters cannot be
teleported. A vehicle that has not arrived when the hour ends counts with
the time it has driven so far.

\begin{table}[htbp]
    \centering
    \caption{Genes of the behaviour calibration. vType parameters apply
    to cars, vans, and trucks; options apply to the whole simulation.}
    \label{tab:behaviour-genes}
    \small
    \begin{tabular}{@{}lp{56mm}rr@{}}
        \toprule
        Gene & Effect & Range & Default \\
        \midrule
        \texttt{tau} & desired time headway (s) & 0.8--1.2 & 1.0 \\
        \texttt{sigma} & driver imperfection & 0.1--0.5 & 0.5 \\
        \texttt{impatience} & willingness to force a gap & 0--1 & 0 \\
        \texttt{jmTimegapMinor} & gap accepted from a minor road (s) & 0.5--1.5 & 1.0 \\
        \texttt{jmIgnoreKeepClearTime} & wait before entering a keep-clear area (s) & off, 1--60 & off \\
        \texttt{lcStrategic} & how early a driver changes lanes for the route & 0.5--5 & 1 \\
        \texttt{lcCooperative} & willingness to let others change in & 0.5--1 & 1 \\
        \texttt{lcAssertive} & acceptance of small gaps when changing lanes & 1--5 & 1 \\
        \texttt{lcLookaheadLeft} & how much earlier a driver plans a change to the left & 1--4 & 2 \\
        \texttt{time-to-impatience} & wait until a driver becomes impatient (s) & 30--300 & 180 \\
        \texttt{ignore-junction-blocker} & wait before passing a blocked junction (s) & off, 1--60 & off \\
        \bottomrule
    \end{tabular}
\end{table}

\paragraph{Hotspots.} The best behaviour is simulated again with a
second random seed. An approach that stops at least five vehicles in
both runs is a hotspot; for each, it is noted whether the approach ends
at a traffic light. Hotspots are not edited further. They are capacity
limits of the modelled network at full traffic, for example a missing
lane or a junction with too little throughput, and are discussed as
limitations (Section~\ref{sec:full-traffic}). The stuck vehicles of
the best run are also split by five-minute interval, by approach, and
by vehicle type against each type's share of the demand.

\paragraph{Traffic counts.} The best run is compared with the hourly
motor-traffic counts of the city of Zürich \citep{zurichcounts}. For
each measurement site, the reference is the median of 07:00--08:00
over the measured Tuesdays to Thursdays of \tbd{2024}, without July,
August, and the Christmas holidays. A counting cross-section is the sum
of its directions and is matched to the nearest SUMO edges within
\(50\,\mathrm{m}\) that carry the same street, one per counted
direction. The agreement is given by the GEH statistic,
\(\mathrm{GEH} = \sqrt{2(M-C)^2/(M+C)}\) with \(M\) the simulated and
\(C\) the counted vehicles per hour; GEH below 5 at 85\,\% of the
cross-sections is the usual acceptance level for an hourly model
\citep{dmrb1996}.

\begin{figure}[htbp]
    \centering
    \begin{adjustbox}{max width=\textwidth, max totalheight=0.78\textheight}
    \begin{tikzpicture}[
        font=\footnotesize,
        node distance=3mm and 6mm,
        box/.style={draw, rounded corners=2pt, align=center, fill=white,
            minimum height=7mm, inner sep=2pt, text width=#1},
        box/.default=34mm,
        data/.style={draw, align=center, fill=gray!12, minimum height=7mm,
            text width=#1, trapezium, trapezium left angle=75,
            trapezium right angle=105, trapezium stretches body},
        data/.default=30mm,
        check/.style={draw, diamond, aspect=2.2, align=center, fill=white,
            inner sep=1pt, text width=22mm},
        osm/.style={box=#1, fill=green!10},
        osm/.default=34mm,
        google/.style={box=#1, fill=blue!8},
        google/.default=34mm,
        arrow/.style={-{Stealth[length=2mm]}, thick},
        tag/.style={font=\scriptsize, fill=white, inner sep=1pt},
    ]
        \node[data] (net) {Unedited SUMO network\\and OSM extract};
        \node[data, right=18mm of net] (plans) {eqasim \texttt{output\_plans}\\and MATSim network};
        \node[box, below=of net] (correct) {\textbf{Filtering}\\rules C1--C5,\\tram access};
        \node[box, below=of plans] (demand) {\textbf{Demand cut}\\follow routes, trim to\\the connected network};
        \node[data, below=of correct] (corrected) {Corrected network};
        \node[data, below=of demand] (scale) {Car, van and truck trips\\07:00--08:00, sample};

        \coordinate (mid) at ($(corrected.south)!0.5!(scale.south)$);
        \node[osm=46mm, below=7mm of mid] (osmstep) {\textbf{OpenStreetMap step}\\connect links that stop where\\OpenStreetMap continues};
        \node[box=52mm, below=of osmstep] (sim) {\textbf{Simulate the National demand}\\vehicles stuck for 120\,s (stage 1)\\or 300\,s (stage 2) are teleported;\\counted from SUMO statistics};
        \node[check, below=of sim] (stuck) {Vehicles stuck?};

        \node[google=27mm, below=9mm of stuck] (turns) {\textbf{Turn genes}\\Google disagrees\\with the network};
        \node[google=27mm, left=4mm of turns] (lanes) {\textbf{Lane genes}\\wrong-lane movement\\Google drives};
        \node[box=27mm, right=4mm of turns] (bans) {\textbf{Ban genes}\\lane-change bans\\near wrong-lane stops};

        \node[box=60mm, below=of turns] (evolve) {\textbf{Evolutionary search}\\8 individuals, 2 elites, uniform crossover,\\one gene flipped, at most 6 generations};
        \node[box=46mm, below=of evolve] (score) {Build network, simulate, score\\stuck vehicles + unroutable trips};
        \node[osm=46mm, below=of score] (review) {\textbf{OpenStreetMap review}\\reject a new link to nowhere\\or a split network};
        \node[check, below=of review] (better) {Better\\network?};

        \node[osm=56mm, below=6mm of better] (final) {\textbf{Final check}\\OpenStreetMap check and simulation};
        \node[check, below=of final] (stage) {Stage 1\\(sample)?};
        \node[box=40mm, left=8mm of stage] (scaleup) {\textbf{Scale to full population}\\\(\times\,1/\)sample share,\\clones \(\pm 30\,\mathrm{min}\), free ends\\moved \(\le 500\,\mathrm{m}\) to homes\\and workplaces};
        \node[box=40mm, anchor=north west] (behaviour) at ($(stage.east |- stage.north)+(8mm,0)$)
            {\textbf{Stage 3a: signal timing}\\Webster from the full\\demand, actuated, max\\green \(1.25\times\); short\\edges with collisions joined};
        \node[google=40mm, below=of behaviour] (calib) {\textbf{Stage 3b: driver behaviour}\\11 genes, 6 individuals,\\2 elites, 3 generations;\\score: teleported \% + not\\inserted \% + travel-time\\bias beyond \(\pm 10\,\%\)};
        \node[box=40mm, below=of calib] (hotspots) {\textbf{Hotspots and counts}\\second seed; approaches\\stuck in both runs; GEH\\against Zürich counts};
        \node[box=40mm, left=8mm of hotspots] (result) {\textbf{Result at full traffic}\\network errors, capacity,\\teleported share, travel\\time against Google, vs.\\the unrepaired network};

        \draw[arrow] (net) -- (correct);
        \draw[arrow] (correct) -- (corrected);
        \draw[arrow] (plans) -- (demand);
        \draw[arrow] (demand) -- (scale);
        \draw[arrow] (corrected.south) |- ($(osmstep.north)+(0,4mm)$) -- (osmstep.north);
        \draw[arrow] (scale.south) |- ($(osmstep.north)+(0,4mm)$) -- (osmstep.north);
        \draw[arrow] (osmstep) -- (sim);
        \draw[arrow] (sim) -- (stuck);
        \draw[arrow] (stuck) -- node[tag, right, pos=0.25] {yes} (turns);
        \draw[arrow] (stuck.south) -- ++(0,-4mm) -| (lanes.north);
        \draw[arrow] (stuck.south) -- ++(0,-4mm) -| (bans.north);
        \draw[arrow] (lanes.south) -- (lanes.south |- evolve.north);
        \draw[arrow] (turns) -- (evolve);
        \draw[arrow] (bans.south) -- (bans.south |- evolve.north);
        \draw[arrow] (evolve) -- (score);
        \draw[arrow] (score) -- (review);
        \draw[arrow] (review) -- (better);
        \draw[arrow] (better) -- node[tag, right, pos=0.8] {no} (final);
        \draw[arrow] (final) -- (stage);
        \draw[arrow] (stage.east) -- node[tag, above] {no} (stage.east -| behaviour.west);
        \draw[arrow] (behaviour) -- (calib);
        \draw[arrow] (calib) -- (hotspots);
        \draw[arrow] (hotspots) -- (result);

        \draw[arrow] (review.east) -- ++(12mm,0) |- node[tag, pos=0.25, rotate=90, anchor=south] {next individual} (evolve.east);
        \draw[arrow] (better.west) -- node[tag, above] {yes: next round} ++(-38mm,0) |- (osmstep.west);
        \draw[arrow] (stuck.east) -- node[tag, above] {no} ++(38mm,0) |- (final.east);
        \draw[arrow, dashed] (stage.west) -- node[tag, above] {yes} (scaleup.east);
        \draw[arrow, dashed] (scaleup.west) -- ++(-4mm,0) |- node[tag, pos=0.25, rotate=90, anchor=south] {stage 2: same rounds at full traffic} ($(osmstep.west)+(0,2mm)$);

        \begin{scope}[on background layer]
            \node[draw=gray, dashed, rounded corners=4pt, inner sep=3mm,
                fit=(osmstep) (sim) (stuck) (lanes) (bans) (evolve) (review) (better),
                label={[font=\scriptsize\itshape, anchor=south west, align=left]south west:one round\\stage 1: up to 10\\stage 2: up to 3}] {};
        \end{scope}
    \end{tikzpicture}
    \end{adjustbox}
    \caption{Filtering, evolutionary repair, and behaviour calibration.
    Stage~1 repairs the network with the sampled National demand;
    stage~2 repeats the rounds with the demand scaled to the full
    population, starting from the repaired network; stage~3 sets the
    signal timing from the full demand and then calibrates driver
    behaviour on that network without changing it. Green steps compare the network with OpenStreetMap, and
    blue steps use Google Routes.}
    \label{fig:repair-flow}
\end{figure}

A network without stuck vehicles in this simulation is not proven
correct on the ground. Each kept edit is supported by Google or by
OpenStreetMap, but Google and OpenStreetMap can disagree with each
other; at the A1/A4 Nordumfahrung, for example, the lane connection
that Google drives is not the one given by the OpenStreetMap
\texttt{destination:lanes} tags. The kept edits are therefore listed
with the results so that they can be checked individually.

% ===== (C) replaces the empty \subsection{Demand Transfer and Trip Retention} =====
\subsection{Demand Transfer and Trip Retention}
\label{subsec:demand-transfer}

The car demand for SUMO is taken from two files of the eqasim
National scenario: the selected plans in \texttt{output\_plans.xml.gz}
and the MATSim network of the same run. Legs with the modes car and truck are used.
Instead of matching the origin and destination of each trip, the
converter follows the MATSim route of each leg link by link. This keeps
traffic that only passes through Zürich, which an origin and
destination filter would lose.

Each MATSim link is converted from the Swiss LV95 coordinates to the
coordinates of the SUMO network. It is assigned to the SUMO passenger
edge nearest to its midpoint, provided that the edge lies within
\(40\,\mathrm{m}\) and runs within \(45^\circ\) of the link's
direction. The time at which the vehicle reaches each link is taken
from the free-flow travel times along the route, stretched so that
they add up to the travel time MATSim recorded for the leg. Where a
route leaves the SUMO network for more than \(300\,\mathrm{s}\) and
returns, it is split into separate pieces. Each piece becomes one SUMO
trip, departing when the vehicle enters the network. A trip is kept
when it enters the network between 07:00 and 08:00. Legs that started
up to three hours earlier are read as well, so that long trips that
reach Zürich within the hour are included.

The ends of each piece are then trimmed to the largest strongly
connected part of the SUMO passenger network, in which every edge can
be reached from every other. Snapping at the boundary of the network
can otherwise place a trip on a short stub with no way in or out. A
piece whose first and last edges are the two directions of one street
only touched the boundary and is dropped, as is a piece that is left
with fewer than two edges. Cars, vans, and trucks are written as
separate vehicle types. The light commercial vehicles of eqasim
(agents \texttt{lcv\_*}, mode truck) are vans of \(6.5\,\mathrm{m}\);
the other trucks are \(12\,\mathrm{m}\) long, with a maximum speed of
\(25\,\mathrm{m/s}\). For the sampled demand every trip enters on
the lane that best suits its route and at the highest speed that is
safe at that point, so that vehicles joining a motorway at the boundary
do not start from a standstill.

The converter reports how many route pieces were kept unchanged,
trimmed, or dropped for each of these reasons.
Table~\ref{tab:demand-retention} lists the counts for the morning
hour.

\begin{table}[htbp]
    \centering
    \caption{Route pieces of the National car and truck demand that
    enter the Zürich network between 07:00 and 08:00.}
    \label{tab:demand-retention}
    \begin{tabular}{lr}
        \toprule
        Outcome & Pieces \\
        \midrule
        Kept unchanged & \tbd{kept} \\
        Trimmed to the connected network & \tbd{trimmed} \\
        Only touches the boundary & \tbd{boundary} \\
        Too short after trimming & \tbd{too short} \\
        Off the main network & \tbd{off} \\
        \midrule
        SUMO trips written & \tbd{trips} \\
        \bottomrule
    \end{tabular}
\end{table}

\subsection{Scaling to the Full Population}
\label{subsec:demand-scaling}

The sampled trips are scaled to the full population with the method of
\citet[Section~3.3.3]{baggenstos2026}. The scale factor is the inverse
of the sample share, which is read from the eqasim run configuration
(\texttt{sampleSize} in \texttt{output\_config.xml}); for a share of
\tbd{0.1} the factor is \tbd{10}.

\paragraph{Trip cloning.} Every original trip is kept unchanged, and
the factor minus one clones are added. A clone's identifier is that of
its original with a clone number appended, so that each clone can be
traced back to its source trip.

\paragraph{Temporal perturbation.} The departure of each clone is
shifted by an offset drawn uniformly from
\(\pm 30\,\mathrm{min}\). \citet{baggenstos2026} keeps a clone when the
shifted time lies in the simulated day. Here only one hour is
simulated, so a clone is kept when it departs between 07:00 and 08:00.
For this, the trips are cut as in Section~\ref{subsec:demand-transfer}
but for 06:30 to 08:30. The originals outside the hour are not
simulated; only their clones can be shifted into it, as clones near
midnight are in the daily run. The expected number of trips in the hour
is then the factor times the sampled trips of the hour.

\paragraph{Spatial perturbation.} The free ends of each clone are moved
within \(500\,\mathrm{m}\) of the original ends. An end is free when it
lies at an activity, that is, when the vehicle is on the SUMO network
within \(60\,\mathrm{s}\) of leaving or reaching its activity; an end
where the route crosses the network boundary stays fixed. As in
Table~4 of \citet{baggenstos2026}, a trip inside the network moves both
ends, an entering trip its end, a leaving trip its start, and a
pass-through trip neither.

\paragraph{Amenity-based relocation.} The new end is drawn from the
home, work, and education locations of the eqasim agents within the
radius, weighted by the number of agents at each location. eqasim
builds these locations from STATPOP and STATENT, the registers
\citet{baggenstos2026} uses directly. Where the circle holds no
location, a uniform point in it is used. The point is placed on the
nearest drivable edge of the connected network within
\(150\,\mathrm{m}\), the distance \citet{baggenstos2026} uses for
internal trips; a point with no edge in reach is drawn again, at most
twenty times, before the original edge is kept. All random draws use a
fixed seed (30), so the scaled demand is reproducible. As in his SUMO
set-up, every scaled trip enters on the best lane at the average speed
of that lane.

% ===== (D) replaces \subsection{Turning movements checked against Google} and \subsection{Remaining gridlocks}, including the table =====
\subsection{Repair with the National Demand}
\label{sec:national-repair}

Stage~1 of the repair (Section~\ref{subsubsec:evolutionary-turns}) was
run with the sampled National demand of \tbd{trips} trips. The
corrected network produced \tbd{stuck} stuck vehicles and \tbd{n}
unroutable trips in the first round. The search ran for \tbd{rounds}
rounds and kept \tbd{n} turn edits, \tbd{n} lane connections, and
lifted \tbd{n} lane-change bans. The OpenStreetMap step connected
\tbd{n} links. The final OpenStreetMap check found
\tbd{no link to nowhere, and the network is connected}.
Table~\ref{tab:repair-edits} lists the kept edits.

\begin{table}[htbp]
    \centering
    \caption{Edits kept by the evolutionary repair.}
    \label{tab:repair-edits}
    \begin{tabular}{llll}
        \toprule
        Edit & Street & From & To \\
        \midrule
        \tbd{add lane} & \tbd{street} & \texttt{\tbd{edge}} & \texttt{\tbd{edge}} \\
        \bottomrule
    \end{tabular}
\end{table}

\subsection{Full-Traffic Check}
\label{sec:full-traffic}

With the demand scaled to the full population, \tbd{trips} trips are
due to enter the network during the hour, of which SUMO inserted
\tbd{inserted}. Table~\ref{tab:full-traffic} compares the unrepaired
network, the network repaired in stage~1, the network after stage~2,
the same network with the Webster signals, and with the calibrated
behaviour of stage~3. A vehicle counts as stuck after
\(300\,\mathrm{s}\) without moving. All numbers are taken from SUMO's
statistic output.

\begin{table}[htbp]
    \centering
    \caption{Stuck vehicles, teleported share, and travel time with the
    National demand scaled to the full population.}
    \label{tab:full-traffic}
    \small
    \begin{tabular}{@{}p{50mm}rrrrr@{}}
        \toprule
        & Unrepaired & Stage 1 & Stage 2 & Signals & Behaviour \\
        \midrule
        Teleported: wrong lane & \tbd{n} & \tbd{n} & \tbd{n} & \tbd{n} & \tbd{n} \\
        Teleported: collision & \tbd{n} & \tbd{n} & \tbd{n} & \tbd{n} & \tbd{n} \\
        Teleported: queue, yield & \tbd{n} & \tbd{n} & \tbd{n} & \tbd{n} & \tbd{n} \\
        Teleported share (\%) & \tbd{p} & \tbd{p} & \tbd{p} & \tbd{p} & \tbd{p} \\
        Not inserted (\% of loaded) & \tbd{p} & \tbd{p} & \tbd{p} & \tbd{p} & \tbd{p} \\
        Unroutable trips & \tbd{n} & \tbd{n} & \tbd{n} & \tbd{n} & \tbd{n} \\
        Travel-time bias against Google (\%) & -- & -- & -- & \tbd{b} & \tbd{b} \\
        Cross-sections with GEH \(<5\) (\%) & -- & -- & -- & -- & \tbd{p} \\
        \bottomrule
    \end{tabular}
\end{table}

The signal variants are compared in Table~\ref{tab:signals}. The
stage~2 network with SUMO's default behaviour and the Webster signals
is the first
individual of stage~3, so its travel-time bias is known as well.
Table~\ref{tab:behaviour-result} lists the calibrated values.

\begin{table}[htbp]
    \centering
    \caption{Signal programs on the stage~2 network with SUMO's default
    behaviour. The actuated Webster programs are used from here on.}
    \label{tab:signals}
    \small
    \begin{tabular}{@{}lrrrr@{}}
        \toprule
        Programs & Teleported (\%) & Not inserted (\%) & Stuck at signals & Time loss (s) \\
        \midrule
        Guessed by netconvert & \tbd{p} & \tbd{p} & \tbd{n} & \tbd{t} \\
        Webster, actuated & \tbd{p} & \tbd{p} & \tbd{n} & \tbd{t} \\
        Webster, fixed time, coordinated & \tbd{p} & \tbd{p} & \tbd{n} & \tbd{t} \\
        \bottomrule
    \end{tabular}
\end{table}

\begin{table}[htbp]
    \centering
    \caption{Calibrated driver behaviour (stage~3), checked with a
    second random seed.}
    \label{tab:behaviour-result}
    \begin{tabular}{lrr}
        \toprule
        Gene & SUMO default & Calibrated \\
        \midrule
        \texttt{tau} & 1.0 & \tbd{v} \\
        \texttt{sigma} & 0.5 & \tbd{v} \\
        \texttt{impatience} & 0 & \tbd{v} \\
        \texttt{jmTimegapMinor} & 1.0 & \tbd{v} \\
        \texttt{jmIgnoreKeepClearTime} (s) & off & \tbd{v} \\
        \texttt{lcStrategic} & 1 & \tbd{v} \\
        \texttt{lcCooperative} & 1 & \tbd{v} \\
        \texttt{lcAssertive} & 1 & \tbd{v} \\
        \texttt{lcLookaheadLeft} & 2 & \tbd{v} \\
        \texttt{time-to-impatience} (s) & 180 & \tbd{v} \\
        \texttt{ignore-junction-blocker} (s) & off & \tbd{v} \\
        \midrule
        Teleported share, seed 23 / seed 42 (\%) & \tbd{p} & \tbd{p} / \tbd{p} \\
        Not inserted, seed 23 / seed 42 (\%) & \tbd{p} & \tbd{p} / \tbd{p} \\
        Mean travel time SUMO / Google (min) & \tbd{t} / \tbd{t} & \tbd{t} / \tbd{t} \\
        Mean relative error, RMSE & \tbd{e} & \tbd{e} \\
        \bottomrule
    \end{tabular}
\end{table}

\begin{table}[htbp]
    \centering
    \caption{Hotspots: approaches that stop at least five vehicles in
    both runs of the calibrated behaviour.}
    \label{tab:gridlocks}
    \begin{tabular}{llcrr}
        \toprule
        Approach and junction & What happened & Signal & Seed 23 & Seed 42 \\
        \midrule
        \tbd{approach / junction} & \tbd{kind} & \tbd{yes} & \tbd{n} & \tbd{n} \\
        \bottomrule
    \end{tabular}
\end{table}

% ===== (E) add under \subsection{Limitations and Sources of Uncertainty} =====
The repair is judged by one simulated morning hour, which starts with
an empty network. Its score counts vehicles that SUMO teleports after
\(120\,\mathrm{s}\) (stage~1) or \(300\,\mathrm{s}\) (stages~2 and~3)
without moving, which depends on SUMO's lane-changing and junction
models as well as on the network. The full traffic is obtained by
cloning the sampled trips. The clones are spread over
\(\pm 30\,\mathrm{min}\) and their free ends over \(500\,\mathrm{m}\),
but they keep the boundary crossings of their originals, so the
through traffic is more concentrated than in reality. MATSim instead
scales road capacity down to the sample, so its congestion and SUMO's
are not directly comparable.

Each edit is supported by Google Routes or OpenStreetMap, but the two
sources disagree at some junctions, and neither has been checked on
site. The behaviour calibration fits one parameter set for all drivers
to sixty Google travel times from one morning, with a single seed in
the search; the calibrated values describe this network and demand,
not drivers in general. The signal programs are not the programs of the
city: their phases are those netconvert guessed, and their timing
follows Webster's method from the simulated demand, so a hotspot at a
traffic light can come from the timing as much as from the junction.
Lane-change bans are respected by the routes chosen before departure;
a vehicle that is rerouted later can still be sent onto a route that
the bans make impossible, except near the edges where this was found.
The hotspots that remain at full
traffic (Table~\ref{tab:gridlocks}) cannot be removed by changing turns,
lanes, signal timing, or behaviour. They are capacity limits of the modelled network,
such as a missing lane or a junction with too little throughput, and
they bound how far the full-traffic results can be read at those
locations.
